\subsection{Переменные действие-угол.}
Не всегда гамильтонову систему хочется интегрировать в квадратурах по Гамильтону-Якоби.
Иногда может интересовать представление в таких координатах, что уравнения Гамильтона будут иметь вид
\[
    \begin{cases}
        \dot{\bt{p}} = \bt{0} \\
        \dot{\bt{q}} = const.
    \end{cases}
\]
В общем случае в этих координатах система движется по $n$-мерному тору. \\
Такие координаты называются <<переменные действие-угол>>.
Покажем как их можно найти в одномерном варианте для системы с гамильтонианом $H = H(q, p)$.
Пусть $H(q, p) = h$ и импульс разрешается относительно $q, h$, то есть $p = p(q, h)$.
Пусть движение системы <<периодично>>, то есть или траектории движения замкнуты (на фазовой плоскости) или существует такой $q_0$, что $q + q_0$ не меняет конфигурацию системы. \\
В таком случае определим действие как $I = \dfrac{1}{2\pi}\oint p(q, h)dq$ по этому самому периоду. \\
Покажем, что замена переменных $(q, p) \mapsto (\phi, I)$ будет канонической.
Для этого найдём производящую функцию. По способу введения, $I = I(h)$.
Тогда $S(q, I) = \int p(q, h(I))dq$ -- искомая. \\
Новым гамильтонианом будет $\tilde{H} = \tilde{H}(I) = h(I)$ и система примет необходимый вид.

\subsection{Классическая теория возмущений.}
Предположим, что гамильтонова система с гамильтонианом $H(\bt{p}, \bt{q})$ представлена в виде:
\[
    H(\bt{p}, \bt{q}) = H_0(\bt{p}, \bt{q}) + \epsilon H_1(\bt{p}, \bt{q}).
\]
где $H_0$ -- гамильтониан интегрируемой системы, $H_1$ -- некоторое <<возмущение>> и $\epsilon \ll 1$.
Будем искать приближенное решение для $H$ в виде суммы точного решения для $H_0$ и малой поправки, связанной с $H_1$.
Будем считать что гамильтониан задан в переменных действие-угол, то есть $H = H(\bt{I}, \bt{\theta})$ и $H_0 = H_0(\bt{I})$.
Тогда уравнения для интегрируемой части имеют вид
\[
    \begin{cases}
        \dot{\bt{\theta}} = w(\bt{I}), \\
        \dot{\bt{I}} = 0.
    \end{cases}
\]
Представим $H_1$ в виде ряда Фурье:
\[
    H_1 = \overline{H}_1(\bt{I}) + \sum_{\bt{k} \in \Z^m} c_k(\bt{I}) e^{i \langle \bt{k}, \bt{\theta} \rangle}.
\]
Наша цель -- найти такое каноническое преобразование, что гамильтониан новой системы не зависит от $\epsilon$ в линейном приближении \[\mathcal{H} = H_0^*(\bt{I}^*) + \epsilon^2 H_2^*(\bt{I}^*, \bt{\theta}^*).\]
Если это получится, то мы получим решение изначальной системы с точностью до линейных поправок по $\epsilon$.
Предположим, что производящая функция этого преобразования имеет вид
\[
    S(\bt{\theta}, \bt{I}^*) = \bt{\theta} \cdot \bt{I}^* + \epsilon S_1^{(1)}(\bt{\theta}, \bt{I}^*).
\]
Поскольку $\bt{\theta}^* = \frac{\partial S}{\partial \bt{I}^*}$ и $\bt{I} = \frac{\partial S}{\partial \bt{\theta}}$, то канонические уравнения примут вид
\[
        \bt{\theta}^* = \bt{\theta} + \epsilon \dfrac{\partial S_1^(1)}{\partial \bt{I}^*}, \quad
        \bt{I}^* = \bt{I} + \epsilon \dfrac{\partial S_1^{(1)}}{\partial \bt{\theta}^*}.
\]
А новый гамильтониан $\mathcal{H}$:
\[
    \mathcal{H} = H_0(\bt{I}^*) + \epsilon \overline{H}_1(\bt{I}^*) + \epsilon \biggr[ \dfrac{\partial H_0}{\partial \bt{I}^*} \cdot \dfrac{\partial S_1^{(1)}}{\partial \bt{\theta}} + \sum\limits_{\bt{k} \in \Z^n} c_k(\bt{I}^*) e^{i\langle \bt{k}, \bt{\theta} \rangle}\biggr].
\]
Теперь для нахождения искомой $S_1^{(1)}$ достаточно проинтегрировать следующее соотношение:
\[
    \dfrac{\partial H_0}{\partial \bt{I}^*} \cdot \dfrac{\partial S_1^{(1)}}{\partial \bt{\theta}} = -\sum\limits_{\bt{k} \in \Z^n} c_k(\bt{I}^*)e^{i\langle \bt{k}, \bt{\theta} \rangle}.
\]
То есть просто формально проинтегрировать ряд Фурье и получить ответ.
Увы, но далеко не факт что это приведёт к решению задачи: для получения решения нужно проверить сходимость ряда.
Если радиус сходимости ряда равен нулю, или ряд в принципе не сходится, то полученное разложение не годится.

В случае успеха, поиск решения можно продолжить и для приближений более высокой степени -- сначала рассматривать разложение возмущения для более высокой степени по $\epsilon$, а потом рассматривать производящую функцию, зависящую от более высоких степени $\epsilon$ и приходить к решению до $o(\epsilon^n)$.
